% Part: normal-modal-logic
% Chapter: completeness
% Section: completeness-K

\documentclass[../../../include/open-logic-section]{subfiles}

\begin{document}

\olfileid{nml}{com}{cmk}

\olsection{\Log{K} चे निर्धारण आणि संपूर्णता}

आता संपूर्णता सिद्ध करण्यासाठी कॅनॉनिकल प्रतिमान
वापरता येईल. $\Sigma$ साठीच्या कॅनॉनिकल
प्रतिमानात सत्य असलेली !!{formula}s म्हणजे
नेमकी $\Sigma$-!!{derivable} सूत्रे असतात;
या तथ्यावरून संपूर्णता मिळते. हा गुणधर्म
असलेले प्रतिमान~$\Sigma$ चे
\emph{निर्धारण करते} असे म्हणतात.

\begin{defn}
  सर्व !!{formula}s पैकी प्रत्येक~$!A$ साठी
  $\mSat{M}{!A}$ असेल तेव्हा आणि तेव्हाच
  $\Sigma \Proves !A$ असेल, तर
  प्रतिमान~$\mModel{M}$ हे सामान्य मोडल
  तर्कशास्त्र~$\Sigma$ चे
  \emph{निर्धारण करते}.
\end{defn}

\begin{thm}[निर्धारण]\ollabel{thm:determination}
   $\mSat{M^\Sigma}{!A}$ असेल तेव्हा
   आणि तेव्हाच $\Sigma \Proves !A$.
\end{thm}

\begin{proof}
  $\mSat{M^\Sigma}{!A}$ असेल, तर प्रत्येक
  संपूर्ण $\Sigma$-सुसंगत $\Delta$ साठी
  $\mSat{M^\Sigma}{!A}[\Delta]$.
  म्हणून सत्यता पूर्वप्रमेयानुसार प्रत्येक
  संपूर्ण $\Sigma$-सुसंगत $\Delta$ साठी
  $!A \in \Delta$. त्यामुळे
  \olref[lin]{cor:provability-characterization}
  वरून ($\Gamma = \emptyset$ घेऊन)
  $\Sigma \Proves !A$.

  उलट $\Sigma \Proves !A$ असेल, तर
  \olref[ccs]{prop:ccs-properties}\olref[ccs]{prop:ccs-closed}
  नुसार प्रत्येक संपूर्ण $\Sigma$-सुसंगत
  $\Delta$ मध्ये~$!A$ असते. म्हणून
  सत्यता पूर्वप्रमेयानुसार प्रत्येक~$\Delta \in
  W^\Sigma$ साठी
  $\mSat{M^\Sigma}{!A}[\Delta]$;
  म्हणजे $\mSat{M^\Sigma}{!A}$.
\end{proof}

\Log{K} साठीचे कॅनॉनिकल प्रतिमान~\Log{K}
चे निर्धारण करत असल्याने
\Log{K} ची संपूर्णता लगेच निष्पन्नतेने मिळते:

\begin{cor}\ollabel{cor:Kcomplete}
  मूलभूत मोडल तर्कशास्त्र~\Log{K}
सर्व प्रतिमानांच्या वर्गाच्या सापेक्ष
संपूर्ण आहे; म्हणजे
$\Entails !A$ असेल, तर $\Log{K}
\Proves !A$.
\end{cor}

\begin{proof}
  प्रतिपरिवर्तनाने: $\Log{K} \Proves/ !A$
  असेल, तर निर्धारणावरून
  $\mSat/{M^{\Log{K}}}{!A}$;
  म्हणून~$!A$ वैध नाही.
\end{proof}

सर्वसाधारणपणे एखादी प्रणाली~$\Sigma$
प्रतिमानांच्या एखाद्या वर्गाच्या सापेक्ष
संपूर्ण आहे हे सिद्ध करण्यासाठी
केवळ निर्धारण पुरेसे नसते. उदाहरणार्थ,
\Log{KTB4} ची परावर्ती, सममित आणि
संक्रमक प्रतिमानांच्या वर्गाच्या सापेक्ष
संपूर्णता विचारात घ्या. प्रणाली~$\Sigma$
साठीचे कॅनॉनिकल प्रतिमान त्या वर्गातील
सदस्य आहे हेही दाखवावे लागते.
कॅनॉनिकल प्रतिमानाच्या रचनेवरून हे
स्पष्टपणे मिळत नाही---काही वेळा
तर ते खरेही नसते!

\end{document}
